% Pillar-5 (MIN-REPORT-RL) evidence exhibits
\section{Evidence: The Stack Moves More Than the Label}
\label{sec:p5-evidence}

\subsection{Exhibit 1: A 17$\times$ Flip from Swapping Only the Backend}
\label{sec:p5-exhibit-backend}
In a backend-swap audit (Stratum B), we held the model, group size $G$, learning
rate, dataset, prompt template, and random seed fixed, and exchanged only the
training backend. Final reward moved from $5.0\%$ on one backend to $84.4\%$ on
the other---a $17\times$ ratio produced by zero algorithmic changes. Nothing in
the run's algorithm label, hyperparameter table, or seed disclosure would let a
reader anticipate this; the difference lives entirely in the substrate
(sampling engine, numerics/precision, loss implementation details, and masking
conventions). For calibration, this single infrastructure swing is an order of
magnitude larger than any label-vs-label gap we measured
(\secref{sec:p5-exhibit-h2h}), and is consistent with the broader observation
that seemingly minor evaluation and serving details shift LLM results by
margins that would be headline findings if attributed to algorithms
\citep{biderman2024lessons, hochlehnert2025sober}.

\subsection{Exhibit 2: The Same Label Denotes Different Experiments}
\label{sec:p5-exhibit-label}
DAPO~\citep{yu2025dapo} is defined by decoupled (asymmetric) clipping
\emph{plus} dynamic sampling: groups whose completions are all-correct or
all-incorrect are resampled, which by construction drives the Zero-Variance
Fraction (the share of groups with identical within-group rewards; see the
companion Pillar-2 paper of \ours{}) toward zero. In our cross-stack audit
(Stratum B), ``DAPO'' run on an open Colab trainer with true dynamic sampling
produced mean ZVF $0.00$, exactly as the algorithm's definition requires. The
same label run on a closed training stack---where dynamic sampling is not
expressible and the practitioner substitutes the standard surrogate of GRPO
with asymmetric clipping---produced mean ZVF $0.58$. Same label, opposite
telemetry: one run never wastes a group, the other spends more than half its
groups on zero-gradient updates. Any downstream comparison citing both runs as
``DAPO'' is comparing two different objects. This is not an indictment of
either stack; it is an indictment of the label as an experimental
specification.

\subsection{Exhibit 3: Four Labels Within Noise on One Fixed Stack}
\label{sec:p5-exhibit-h2h}
Conversely, when the stack \emph{is} held fixed, labels stop separating. In a
12-cell head-to-head on the Tinker stack (Qwen3.5-4B, GSM8K, three seeds per
arm; Stratum B), the four labeled arms landed within noise of one another
(\tableref{tab:p5-h2h}): the full spread in mean last-10 training reward is
$0.034$ ($0.710$--$0.744$), comparable to seed-level variation under the
\ours{} seed policy (\secref{sec:setup}). Mean ZVF varies somewhat more
($0.500$--$0.578$) but with no arm approaching the ZVF~$0.00$ that true dynamic
sampling produces (Exhibit 2)---consistent with all four arms being, at this
stack's level of description, minor variations of one group-relative loss. The
juxtaposition of Exhibits 1--3 is the core of our position: a factor the field
reports meticulously (the label) moved outcomes by $\leq 0.034$ reward on a
fixed stack, while a factor the field routinely omits (the backend) moved
outcomes by $0.794$.

\begin{table}[t]
\centering
\caption{12-cell head-to-head on one fixed stack (Tinker; Qwen3.5-4B; GSM8K;
3 seeds per arm; internal program records, June 2026). Arms are within noise
of one another on final reward; no arm reaches the ZVF regime that the same
labels produce on stacks where their defining mechanisms are actually
expressible.}
\label{tab:p5-h2h}
\small
\begin{tabular}{@{}lcc@{}}
\toprule
Arm (label as configured) & Mean last-10 reward & Mean ZVF \\
\midrule
dapo (asymmetric-clip surrogate) & 0.744 & 0.578 \\
drgrpo                           & 0.742 & 0.500 \\
grpo                             & 0.723 & 0.567 \\
gspo                             & 0.710 & 0.511 \\
\bottomrule
\end{tabular}
\end{table}

\subsection{Exhibit 4: Cross-Library Variation in the Benchmark Itself}
\label{sec:p5-exhibit-library}
The pattern is not unique to our audits. Within \ours{} (Stratum A), the
companion Pillar-2 paper documents per-library ZVF aggregations across the
GRPO-family implementations of the benchmark roster
(\tableref{tab:implementations}), with collapse, drift, and plateau rates that
vary by library under matched task and model settings; the released
matched-configuration dumps (one exact hyperparameter dump per framework for
the canonical Qwen3-8B\,+\,GSM8K run) exist precisely because ``the same
configuration'' across frameworks required a nontrivial mapping exercise to
even define. The same-stack PPO-vs-GRPO isolation runs in the benchmark make
the complementary point: once the stack is pinned, algorithm effects can be
measured cleanly---which is the outcome our reporting standard is designed to
make routine rather than exceptional.

\subsection{Exhibit 5: Quantitative Decomposition---Algorithm vs Stack}
\label{sec:p5-exhibit-eta2}
Exhibit 3 shows the qualitative version: four labels within noise on one
fixed stack. To make the comparison quantitatively testable, we decompose
variance into $\eta^{2} = \mathrm{SS}_{\text{between}} / \mathrm{SS}_{\text{total}}$
across two axes on the live N2 corpus
(\texttt{experiments/results/n2\_reward\_tensor\_resume/n2\_metrics.tsv} and
\texttt{experiments/results/groupsize\_zvf\_sweep.tsv}; analysed by
\texttt{scripts/p5p8/stack\_eta2.py}, see
\texttt{experiments/results/p5p8/stack\_eta2.tsv}):

\begin{table}[t]
\centering
\caption{Variance-explained by axis on the live N2 corpus. The algorithm
axis is the four-method same-stack N2 run (GRPO / AERO / GIFT / AREAL,
$n=40$ steps each, one seed); the group-size axis is the four-$G$ sweep on
the Tinker stack (3 seeds per arm). The algorithm axis explains $\le 6.3\%$
of telemetry variance; the group-size axis explains $\ge 22\times$ more for
ZVF and $\ge 133\times$ more for reward. Loss is excluded because it is
method-defined (each method has its own loss surface).}
\label{tab:p5-eta2}
\small
\begin{tabular}{@{}lccc@{}}
\toprule
Telemetry & $\eta^{2}$(algorithm) & $\eta^{2}$($G$) & ratio $G$/alg \\
\midrule
ZVF       & 0.0454 & 1.0000 & 22$\times$  \\
PCD       & 0.0357 & ---    & ---    \\
reward\_mean & 0.0075 & 1.0000 & 133$\times$ \\
mean\_len  & 0.0631 & ---    & ---    \\
cv\_len    & 0.0457 & ---    & ---    \\
\bottomrule
\end{tabular}
\end{table}

The structural reading: when the stack is held fixed, swapping GRPO for
AERO / GIFT / AREAL moves ZVF by $\le 6.3\%$ of the variance budget;
changing $G$ moves it by an order of magnitude more. This is the
quantitative counterpart to Exhibit~3 and the empirical content of the
``Report the Stack, Not the Label'' thesis: the algorithm axis is
under-identified against the stack axis on this corpus. (Caveat: the
algorithm-axis $\eta^{2}$ is from a single seed; the \emph{direction}
$\eta^{2}(\text{algorithm}) \ll \eta^{2}(G)$ is robust, the magnitude
is not.)

\subsection{Exhibit 6: MIN-REPORT Field Coverage at $n=98$ Manifests}
\label{sec:p5-exhibit-coverage}
To stress-test the standard at scale, we audited every manifest in
\texttt{experiments/results/mega\_20260704/manifests/} ($n=98$ cells) for
the seven MIN-REPORT items of \secref{sec:p5-stack} using
\texttt{scripts/p5p8/minreport\_coverage.py}; full table at
\texttt{experiments/results/p5p8/minreport\_field\_coverage.tsv}:

\begin{table}[t]
\centering
\caption{Per-item coverage of the seven-item MIN-REPORT standard against
$n=98$ live mega-campaign manifests. All seven items are 100\% present
\emph{as keys}, but Item 2 (Reference policy \& KL) reports a concrete
\texttt{kl-*} value in $0/98$ ($0.0\%$)---the corpus is uniformly
sampling-only with no KL regulariser, but the value \texttt{"n/a"} is
indistinguishable from an emission bug.}
\label{tab:p5-coverage}
\small
\begin{tabular}{@{}clcc@{}}
\toprule
\# & Item & present & validated \\
\midrule
1 & Loss form              & 100.0\% & 100.0\% \\
2 & Reference policy \& KL & 100.0\% &   0.0\% \\
3 & Sampler / backend / precision & 100.0\% & 100.0\% \\
4 & Per-step ZVF/GU trajectory   & 100.0\% & 100.0\% \\
5 & Group-size schedule          & 100.0\% & 100.0\% \\
6 & Held-out split               & 100.0\% & 100.0\% \\
7 & Decontamination \& parser probe & 100.0\% & 100.0\% \\
\bottomrule
\end{tabular}
\end{table}

All twelve sub-fields enumerated in the standard (ratio level, clip range,
advantage normalisation, dynamic sampling, token mask, KL snapshot, KL
coefficient, KL estimator, precision, decoding parameters, contamination
check, parser probe) score $0/98$ in the live corpus. The standard is
therefore satisfiable as a shallow key set, but not yet enforceable as a
structured stack record---the gap the MIN-REPORT-RL Auditor of
\secref{sec:p5-toolchain} is designed to close.

\subsection{Exhibit 7: Stack-axis $\eta^{2}$ on the 98-cell Live Mega Corpus}
\label{sec:p5-exhibit-mega-eta2}
Exhibit~5 used the 4-method same-stack N2 run to compare the algorithm axis
against the group-size axis ($n=160$ per-step observations, one seed). To test
whether the same direction holds at corpus scale---and across the full
multifactor stack---we ran a one-way $\eta^{2}$ decomposition on the live
mega-campaign ledger (\texttt{experiments/results/mega\_20260704/cells.tsv};
$n=98$ cells = $2$ models $\times$ $3$ task slices $\times$ $5$ $G$ values
$\times$ $2$ temperatures $\times$ $2$ seeds). The five axes analysed are
four \emph{stack} axes (model\_family, task\_slice, $G$, temperature) and one
\emph{noise} axis (seed). Analysis script:
\texttt{scripts/p5p8/mega\_eta2.py}; full output at
\texttt{experiments/results/p5p8/mega\_eta2.tsv}.

\begin{table}[t]
\centering
\caption{Per-axis $\eta^{2}$ on the 98-cell mega corpus. Bold = $\eta^{2}
\ge 0.20$ (DOMINANT verdict). Stack axes dominate every telemetry channel;
seed contributes $\le 0.15\%$ in all channels.}
\label{tab:p5-mega-eta2}
\small
\begin{tabular}{@{}lccccc@{}}
\toprule
axis & \texttt{reward} & \texttt{zvf} & \texttt{pcd} & \texttt{len} & \texttt{std\_len} \\
\midrule
model\_family  & \textbf{0.453} & 0.005 & 0.093 & \textbf{0.301} & 0.159 \\
task\_slice    & \textbf{0.273} & \textbf{0.469} & \textbf{0.451} & \textbf{0.210} & \textbf{0.414} \\
$G$            & 0.030          & \textbf{0.444} & \textbf{0.230} & 0.017 & 0.034 \\
temperature    & 0.000          & 0.009 & 0.009 & \textbf{0.207} & 0.123 \\
seed           & 0.000          & 0.000 & 0.001 & 0.000 & 0.002 \\
\midrule
\emph{stack} $\eta^{2}$ sum & \textbf{0.756} & \textbf{0.927} & \textbf{0.783} & \textbf{0.734} & \textbf{0.729} \\
\bottomrule
\end{tabular}
\end{table}

The four stack axes jointly explain $73\%$--$93\%$ of the variance in every
telemetry channel; the seed axis explains $0.0$--$0.15\%$. The
stack-vs-seed $\eta^{2}$ ratio ranges from $503\times$ (\texttt{std\_len})
to $96{,}128\times$ (\texttt{reward}) and is effectively infinite for
\texttt{zvf} and \texttt{reward} (seed $\eta^{2}=0$ exactly on the 98 cells).

The per-task $G$-axis decomposition sharpens the picture: $\eta^{2}(G)$
for \texttt{zvf} is $0.89$ on \texttt{gsm8k\_easy}, $0.64$ on
\texttt{gsm8k\_hard}, and $0.00$ on \texttt{humaneval\_subset}---the latter
is the all-wrong degenerate regime where the policy fails on every
completion and no $G$ can recover within-group contrast. This recovers
and quantifies the Pillar-2 controller-headroom scope: $G$ is a real lever
on informative tasks ($0.64$--$0.89$) and structurally inert on degenerate
ones ($0.00$). The corpus-scale numbers also refine Exhibit~5: the
algorithm axis was $\le 6.3\%$ on N2 (4 methods, same stack), and on the
98-cell mega corpus the algorithm is \emph{not even a fixed axis}---it
is absorbed into the model\_family or task\_slice cells. What matters for
reported outcomes is which model and which task, not which loss form.

The structural consequence for MIN-REPORT is that the seven items of
\secref{sec:p5-stack} currently span only two of the five axes that
explain the data (Item~5 captures $G$; Item~6 captures the held-out split
but not the task slice). A natural extension---recommended for future
work---is to add an explicit \texttt{model\_family} and
\texttt{task\_slice} declaration alongside the existing seven items, so
that the audit of \secref{sec:p5-exhibit-coverage} can detect the two
axes that explain the most variance in \texttt{reward} and \texttt{len}.

\subsection{Exhibit 8: The MIN-REPORT-RL Auditor Prototype at $n=103$}
\label{sec:p5-exhibit-auditor}
To close the gap between ``we ask for seven items'' and ``we can verify
that seven items were reported,'' we prototyped the
\textsc{Min-Report-RL Auditor} of \secref{sec:p5-toolchain} and ran it
across every manifest in the worktree
($n=103$ = $98$ mega-campaign manifests $+5$ quick $2026$-$07$-$04$
manifests; script \texttt{scripts/p5p8/minreport\_auditor.py}, full
output at
\texttt{experiments/results/p5p8/minreport\_audit.tsv} and
\texttt{\_summary.json}). Each manifest is scored on a weighted $0$--$100$
badge---items~3, 4, 7 each worth $20$ pts (the least-reported and
highest-leverage rows of \tableref{tab:p5-minreport}); items~1, 2, 5, 6
each worth $10$ pts---with an \emph{honest} ``n/a'' declaration (a
recognised value like \texttt{n/a-sampling} for Item~1 or \texttt{n/a}
for Item~2) earning $50\%$ of the item's weight, vs.\ $100\%$ for a
fully-validated value and $0\%$ for a missing key.

\begin{table}[t]
\centering
\caption{MIN-REPORT-RL Auditor per-item coverage on $n=103$ manifests.
Items 3 (backend) and 4 (ZVF trajectory) are the highest-leverage
rows of the standard; both score $\le 64\%$ because the corpus is
closed-stack (Item 3) and the trajectory paths in the manifests
reference files that exist on the live machine but not on the audit
host (Item 4). Item~7 (decontamination + parser probe) drops because
the parser-probe sub-field is reported on $0/103$ manifests---the
exact gap the audit was designed to surface.}
\label{tab:p5-auditor-peritem}
\small
\begin{tabular}{@{}clcc@{}}
\toprule
\# & Item & weight & \% achieved \\
\midrule
1 & Loss form                          & 10  & 24.4\% \\
2 & Reference policy \& KL             & 10  & 24.4\% \\
3 & Sampler / backend / precision      & 20  & 64.0\% \\
4 & Per-step ZVF/GU trajectory         & 20  & 49.5\% \\
5 & Group-size schedule                & 10  & 74.3\% \\
6 & Held-out split                     & 10  & 76.2\% \\
7 & Decontamination \& parser probe    & 20  & 61.8\% \\
\midrule
   & \emph{weighted total}              & 100 & \textbf{55.0} \\
\bottomrule
\end{tabular}
\end{table}

The auditor discriminates \emph{within} the corpus. The badge spans
$[35.8, 56.7]$ with std~$3.1$; the five \texttt{quick\_20260704}
manifests (which use verbose keys like
\texttt{ref\_policy\_kl\_handling} and free-text
\texttt{sampler\_backend\_precision} descriptions) score a mean of
$45.0$ vs.\ $55.5$ for the $98$ compact-key mega manifests. Stratifying
by the five axes of \secref{sec:p5-exhibit-mega-eta2}---the axes
iter-5's $\eta^{2}$ showed explain $73$--$93\%$ of outcome
variance---gives a near-flat per-axis distribution
(\tableref{tab:p5-auditor-stratified}), which is itself the
informative result: when the corpus is a single-stack harvest, no
auditor can differentiate cells \emph{along the stack axes}---the
audit's discriminating power is across corpora, not within one. The
$103$-manifest run is therefore a demonstration that the standard is
\emph{enforceable}: every cell received a verdict, every key was read
or declared missing, and the per-item table surfaces the exact gaps
the iter-1 coverage audit flagged (Item 2 KL, Item 7 parser probe).

\begin{table}[t]
\centering
\caption{MIN-REPORT-RL Auditor mean badge stratified by the five
mega-corpora stack axes. Within a single-stack harvest the auditor
is by construction flat across axes; the discriminating axis is
\emph{corpus family} ($55.5$ vs.\ $45.0$, $\Delta = 10.5$).}
\label{tab:p5-auditor-stratified}
\small
\begin{tabular}{@{}lcc@{}}
\toprule
axis & $n$ & mean badge \\
\midrule
\textit{by corpus}    &             & \\
\quad mega (compact keys)        & 98 & 55.5 \\
\quad quick (verbose keys)        &  5 & 45.0 \\
\textit{by task\_slice} &             & \\
\quad \texttt{gsm8k\_easy}        & 32 & 56.7 \\
\quad \texttt{gsm8k\_hard}        & 32 & 56.7 \\
\quad \texttt{humaneval\_subset} & 32 & 53.3 \\
\textit{by $G$}        &             & \\
\quad $G=2$                       & 20 & 55.6 \\
\quad $G=32$                      & 20 & 56.3 \\
\quad $G=4$                       & 20 & 55.3 \\
\quad $G=16$                      & 20 & 55.2 \\
\quad $G=8$                       & 18 & 55.2 \\
\bottomrule
\end{tabular}
\end{table}

The artifact is at
\texttt{experiments/results/p5p8/figures/minreport\_badge\_dist.png}
(histogram with the four tier thresholds marked: bronze $\ge 50$,
silver $\ge 75$, gold $\ge 90$, wood $\ge 25$, fail below). On the
live corpus the auditor awards $0$ gold, $0$ silver, $99$ bronze,
$4$ wood, $0$ fail---the corpus satisfies the standard as
\emph{key-set completeness} (every key is present) but not as
\emph{declaration completeness} (Item 2's concrete KL value, Item 7's
parser-probe sub-field, Item 4's on-disk trajectory). The fix is
mechanical: extend the manifest emitter with the three missing
declarations, and the badge will move into the silver tier on the
same $103$ cells. This closes the loop the iter-1 audit identified
\emph{operationally}: the auditor prototype does not just count keys,
it scores them, and the per-item table becomes the work-list for the
next iteration of the manifest emitter.
